\section{Precomputed Canonical String Cache}
\label{sec:precomputed}

\subsection{Motivation}

The canonicalization step (especially \texttt{exhaustive} and \texttt{pruned-exhaustive})
is the single most CPU-intensive operation in IsalSR. Its output is a \textbf{pure function}
of the DAG structure --- it depends on no random seed, no training data, and no SR method.
Therefore, for a given configuration $(m, \mathcal{O}, k_{\max})$ (number of variables,
operator set, maximum internal nodes), the mapping
\[
  D \;\mapsto\; \bigl\{\,\text{algorithm} : w_D^{\text{algo}}\,\bigr\}
\]
is deterministic and reusable. Precomputing this mapping eliminates redundant computation
across seeds ($O(30)$), SR methods ($O(3)$), and experiment reruns, yielding an estimated
$90\times$ total CPU time reduction.

\subsection{Cache Data Structure}

The cache is stored as HDF5 files (via \texttt{h5py}) with a JSON sidecar for
human-readable metadata. Each file covers one $(m, \mathcal{O})$ configuration.

\subsubsection{HDF5 Schema}

\begin{verbatim}
cache_{benchmark}_{m}vars_{ops_hash}.h5
|-- attrs/                 (metadata)
|   |-- num_variables      int
|   |-- operator_set       JSON string
|   |-- creation_timestamp ISO 8601
|   |-- git_hash           str
|   `-- total_entries      int
|-- strings/               (all D2S outputs)
|   |-- raw                [N] vlen-str
|   |-- greedy_single      [N] vlen-str
|   |-- greedy_min         [N] vlen-str
|   |-- pruned             [N] vlen-str
|   |-- exhaustive         [N] vlen-str
|   `-- is_canonical       [N] bool
|-- dag_properties/
|   |-- n_nodes            [N] int32
|   |-- n_internal         [N] int32
|   |-- n_edges            [N] int32
|   |-- depth              [N] int32
|   `-- n_var_nodes        [N] int32
|-- timings/               (wall-clock seconds)
|   |-- greedy_single      [N] float64
|   |-- greedy_min         [N] float64
|   |-- pruned             [N] float64
|   `-- exhaustive         [N] float64
|-- correctness/           (cross-validation)
|   |-- exhaustive_timed_out       [N] bool
|   |-- exhaustive_eq_pruned       [N] bool
|   |-- greedy_single_eq_exhaustive [N] bool
|   `-- greedy_min_eq_exhaustive    [N] bool
`-- canonical_index/       (deduplication)
    |-- unique_canonical_pruned    [U] vlen-str
    |-- canonical_to_first         [U] int64
    `-- multiplicity               [U] int64
\end{verbatim}

Where $N$ = total cache entries, $U$ = unique canonical strings ($U \leq N$).

\subsubsection{Fields Stored for Downstream Analysis}

\begin{enumerate}[nosep]
  \item \textbf{Algorithm comparison} (\texttt{strings/}): All four D2S outputs
    stored side-by-side enable offline verification that the exhaustive and pruned
    canonicals agree (Theorem~\ref{thm:pruned_optimality}) and measurement of
    how often greedy variants coincide with the canonical.

  \item \textbf{Timing profiles} (\texttt{timings/}): Per-algorithm wall-clock
    seconds enable cost modeling for the pruning speedup and identification of
    pathological DAGs that trigger exponential backtracking.

  \item \textbf{Search space reduction} (\texttt{canonical\_index/multiplicity}):
    The multiplicity of each unique canonical directly measures the $O(k!)$
    reduction factor --- the ratio $N/U$ is the empirical mean reduction factor.

  \item \textbf{DAG complexity distribution} (\texttt{dag\_properties/}):
    Node count, edge count, and depth characterize the DAG population, enabling
    stratified analysis by expression complexity.

  \item \textbf{Correctness audit} (\texttt{correctness/}): Boolean flags for
    cross-algorithm agreement serve as regression tests: any entry with
    \texttt{exhaustive\_eq\_pruned = False} after the label-aware pruning fix
    (Section~\ref{sec:results}) would indicate a bug.
\end{enumerate}

\subsection{Generation Strategy}

\subsubsection{Sampled Mode (Phase 2)}

For each configuration, generate $N$ random IsalSR strings via
\texttt{random\_isalsr\_string()}, decode each via S2D, and compute all four
D2S algorithms with timing. Invalid strings (VAR-only DAGs, multi-sink DAGs)
are silently skipped. The exhaustive canonical is computed with a 60-second
timeout; if exceeded, the pruned canonical is used as fallback.

Generation is parallelized via SLURM array jobs: each task produces a
\emph{shard} (\texttt{cache\_shard\_\{run\_id\}.h5}), and a follow-up
merge job combines shards with deduplication by raw string.

\subsubsection{CONST Normalization}

All canonical string computations are performed \emph{after} CONST creation
edge normalization (Section~\ref{sec:results}). This ensures that the cache
entries reflect the reduced search space where CONST creation source is
canonicalized to $x_1$.

\subsection{Picasso HPC Deployment}

The cache generation is configured in \texttt{slurm/config.yaml}:

\begin{center}
\begin{tabular}{lccccr}
\toprule
Config & $m$ & $|\mathcal{O}|$ & Strings/task & Shards & Total strings \\
\midrule
Nguyen 1-var & 1 & 8 & 50{,}000 & 20 & 1{,}000{,}000 \\
Nguyen 2-var & 2 & 8 & 50{,}000 & 20 & 1{,}000{,}000 \\
Feynman 1-var & 1 & 10 & 25{,}000 & 20 & 500{,}000 \\
Feynman 2-var & 2 & 10 & 25{,}000 & 20 & 500{,}000 \\
Feynman 3-var & 3 & 10 & 25{,}000 & 20 & 500{,}000 \\
\bottomrule
\end{tabular}
\end{center}

Each shard runs independently with 16\,GB RAM and an 8-hour time limit.
The merge step runs as a single job with 32\,GB RAM.

\subsection{Usage in Experiments}

During SR experiments, the \texttt{CacheManager} provides $O(1)$ lookup
of canonical strings for previously encountered DAGs:

\begin{verbatim}
cache = CacheManager(num_variables, operator_set)
cache.load_hdf5("precomputed/nguyen/nguyen_1var.h5")

entry = cache.lookup(raw_string)  # O(1) dict lookup
if entry is not None:
    canonical = entry.pruned      # or entry.exhaustive
\end{verbatim}

The \texttt{cache.stats} property provides hit/miss statistics and an
estimate of total CPU time saved by precomputation.
